% !PRINTED 5
\weberpage{5}
%<<CONT
and therein $hmn'>hm'n$. One can always take $h$ so large that
further numbers lie between $hmn'$ and $hm'n$, and if $p$ is such a
number, then $p:hnn'$ lies between $\mu$ and $\mu'$.

The points of a straight line can also be regarded as an ordered set
if by greater and smaller one understands some indication of position,
for instance farther right and farther left, or higher and lower.

\begin{center}
\rule{0.18\linewidth}{0.4pt}
\end{center}

A division of an ordered set $\mathfrak M$ into two parts $A,B$ of
such a kind that every element $a$ of $A$ is smaller than every
element $b$ of $B$ is called a \gesperrt{cut} in $\mathfrak M$ and is
suitably denoted by $(A,B)$. Such a cut arises when one selects an
element $\mu$ of $\mathfrak M$, assigns all smaller elements to $A$,
all larger elements to $B$, and assigns $\mu$ itself at will to $A$
or to $B$. Strictly speaking, according as one does the one or the
other, two cuts arise, but we shall always regard them as equal. If,
in a cut $(A,B)$, either $A$ contains a greatest element or $B$ a
least element $\mu$, then we say that $\mu$ produces the cut $(A,B)$.
But it may also occur that neither $A$ has a greatest nor $B$ a least
element.

If every cut in a dense set is produced by a definite element $\mu$,
then the set is called \gesperrt{continuous}.

Continuity as well as density are properties which, by the nature of
the matter, are inaccessible to our sense perception; consequently
they also cannot be demonstrated with rigor in things of the external
world, in spatial magnitudes, intervals of time, or masses, however
much they may seem to lie in the essence of our intuition. But pure
conceptual systems can very well be constructed to which density
without continuity, or also density and continuity, belong\footnotemark[1].
\footnotetext[1]{This was proved by \gesperrt{Dedekind}, to whom we owe
in general the definition of continuity given above. Compare the
writings of \gesperrt{Dedekind}, \emph{Stetigkeit und irrationale
Zahlen}, Brunswick 1872, 1892, and \emph{Was sind und was sollen die
Zahlen?}, Brunswick 1888, 1893. Other}

The \gesperrt{rational fractions} furnish an example of a dense set.
This manifold, which we shall denote by $\mathfrak R$,
%CONT>>
